\section{Customized  Fast Solvers for Diffusion ODEs }
\label{sec:solver}
As highlighted in Sec.~\ref{sec:bg_ode}, discretizing SDEs is generally difficult in high dimensions~\citep[Chap. 11]{kloeden1992Numerical} and it is hard to converge within few steps. In contrast, ODEs are easier to solve, yielding a potential for fast samplers.
However, as mentioned in Sec.~\ref{sec:bg_ode}, the general black-box ODE solver used in previous work~\citep{song2020score} empirically fails to converge in few steps. This motivates us to design a dedicated solver for diffusion ODEs to enable fast and high-quality few-step sampling. 
We start with a detailed investigation of the specific structure of diffusion ODEs.




\subsection{Simplified Formulation of Exact Solutions of Diffusion ODEs}
\label{sec:exact_solution}
The key insight of this work is that given an initial value $\x_s$ at time $s>0$, the solution $\x_t$ at each time $t < s$ of diffusion ODEs in Eq.~\eqref{eqn:diffusion_ode} can be simplified into a very special exact formulation which can be efficiently approximated. 



Our first key observation is that a part of the solution $\x_t$ can be exactly computed by considering the particular structure of diffusion ODEs.
The r.h.s. of diffusion ODEs in Eq.~\eqref{eqn:diffusion_ode} consists of two parts: the part $f(t)\x_t$ is a linear function of $\x_t$, and the other part $\frac{g^2(t)}{2\sigma_t} \epsilonv_\theta(\x_t,t)$ is generally a nonlinear function of $\x_t$ because of the neural network $\epsilonv_\theta(\x_t,t)$. This type of ODE is referred to as \textit{semi-linear} ODE. The black-box ODE solvers adopted by previous work~\cite{song2020score} are ignorant of this semi-linear structure as they take the whole $\hv_\theta(\x_t,t)$ in Eq.~\eqref{eqn:diffusion_ode} as the input, which causes discretization errors of both the linear and nonlinear term. 
We note that for semi-linear ODEs, the solution at time $t$ can be exactly formulated by the \textit{``variation of constants''} formula~\citep{atkinson2011numerical}:
\begin{equation}
\label{eqn:variation_of_constants}
    \x_t = e^{\int_s^t f(\tau)\dd\tau}\x_s
        + \int_s^t \left(e^{\int_\tau^t f(r)\dd r}\frac{g^2(\tau)}{2\sigma_\tau} \epsilonv_\theta(\x_\tau,\tau)\right)\dm\tau.
\end{equation}
This formulation decouples the linear part and the nonlinear part. In contrast to black-box ODE solvers, the linear part is now exactly computed, which eliminates the approximation error of the linear term. However, the integral of the nonlinear part is still complicated because it couples the coefficients about the noise schedule (i.e., $f(\tau),g(\tau),\sigma_\tau$) and the complex neural network $\epsilonv_\theta$, which is still hard to approximate.

Our second key observation is that the integral of the nonlinear part can be greatly simplified by introducing a special variable. Let $\lambda_t \coloneqq \log(\alpha_t / \sigma_t)$ (one half of the log-SNR), then $\lambda_t$ is a strictly decreasing function of $t$ (due to the definition of DPMs as discussed in Sec.~\ref{sec:diffusion-sde}). We can rewrite $g(t)$ in Eq.~\eqref{eqn:ft_gt} as
\begin{equation}
\begin{aligned}
    g^2(t) = \frac{\dd\sigma_t^2}{\dt} - 2\frac{\dd\log\alpha_t}{\dt}\sigma_t^2
    = 2\sigma_t^2\left(\frac{\dd\log\sigma_t}{\dt} - \frac{\dd\log\alpha_t}{\dt}\right) = -2\sigma_t^2\frac{\dd\lambda_t}{\dt}.
\end{aligned}
\end{equation}
Combining with $f(t)=\dd \log \alpha_t / \dt$ in Eq.~\eqref{eqn:ft_gt}, we can rewrite Eq.~\eqref{eqn:variation_of_constants} as
\begin{align}
    \x_t = \frac{\alpha_t}{\alpha_s}\x_s - \alpha_t \int_s^t \left(\frac{\dd\lambda_\tau}{\dd\tau}\right) \frac{\sigma_\tau}{\alpha_\tau} \epsilonv_\theta(\x_\tau,\tau)\dd\tau.
    \label{eqn:analytical_solution_of_time}
\end{align}
As $\lambda(t)=\lambda_t$ is a strictly decreasing function of $t$, it has an inverse function $t_\lambda(\cdot)$ satisfying $t=t_\lambda(\lambda(t))$. We further change the subscripts of $\x$ and $\epsilonv_\theta$ from $t$ to $\lambda$ and denote $\hat\x_{\lambda}\coloneqq \x_{t_\lambda(\lambda)}$, $\hat\epsilonv_\theta(\hat\x_{\lambda},\lambda)\coloneqq \epsilonv_\theta(\x_{t_\lambda(\lambda)}, t_\lambda(\lambda))$. Rewrite Eq.~\eqref{eqn:analytical_solution_of_time} by \textit{``change-of-variable''} for $\lambda$, then we have:



\begin{proposition}[Exact solution of diffusion ODEs]
\label{proposition:exact_solution}
Given an initial value $\x_s$ at time $s>0$, the solution $\x_t$ at time $t\in[0,s]$ of diffusion ODEs in Eq.~\eqref{eqn:diffusion_ode} is:
\begin{equation}
\label{eqn:analytic_solution}
    \x_t = \frac{\alpha_t}{\alpha_s}\x_s - \alpha_t \int_{\lambda_s}^{\lambda_t} e^{-\lambda} \hat\epsilonv_\theta(\hat\x_\lambda,\lambda)\dd\lambda.
\end{equation}
\end{proposition}
We call the integral $\int e^{-\lambda} \hat\epsilonv_\theta(\hat\x_\lambda,\lambda)\dd\lambda$ the \textit{exponentially weighted integral} of $\hat\epsilonv_\theta$, which is very special and highly related to the \textit{exponential integrators} in the literature of ODE solvers~\citep{hochbruck2010exponential}.
To the best of our knowledge, such formulation has not been revealed in prior work of diffusion models.


Eq.~\eqref{eqn:analytic_solution} provides a new perspective for approximating the solutions of diffusion ODEs.
Specifically, given $\x_s$ at time $s$, According to Eq.~\eqref{eqn:analytic_solution}, approximating the solution at time $t$ is equivalent to directly approximating the exponentially weighted integral of $\hat\epsilonv_\theta$ from $\lambda_s$ to $\lambda_t$, which avoids the error of the linear terms and is well-studied in the literature of exponential integrators~\citep{hochbruck2010exponential,hochbruck2005explicit}. Based on this insight, we propose fast solvers for diffusion ODEs, as detailed in the following sections.





\subsection{High-Order Solvers for Diffusion ODEs}
\label{sec:high-order}
In this section, we propose high-order solvers for diffusion ODEs with convergence order guarantee by leveraging our proposed solution formulation Eq.~\eqref{eqn:analytic_solution}.~
The proposed solvers and analysis are highly motivated by the methods of exponential integrators~\citep{hochbruck2010exponential,hochbruck2005explicit} in the ODE literature.


Specifically, given an initial value $\x_T$ at time $T$ and $M+1$ time steps $\{t_i\}_{i=0}^{M}$ decreasing from $t_0=T$ to $t_{M}=0$. Let $\tilde\x_{t_0}=\x_T$ be the initial value. 
The proposed solvers use $M$ steps to iteratively compute a sequence $\{\tilde\x_{t_i}\}_{i=0}^{M}$
to approximate the true solutions at time steps $\{ t_{i} \}_{i=0}^M$. In particular, the last iterate $\tilde\x_{t_M}$ approximates the true solution at time $0$.

In order to reduce the approximation error between $\tilde\x_{t_{M}}$ and the true solution at time $0$, we need to reduce the approximation error for each $\tilde\x_{t_i}$ at every step~\citep{atkinson2011numerical}. Starting with the previous value $\tilde\x_{t_{i-1}}$ at time $t_{i-1}$, according to Eq.~\eqref{eqn:analytic_solution}, the exact solution $\x_{t_{i-1} \to t_i} $ at time $t_{i}$ is given by
\begin{equation}
\label{eqn:analytic_solution_each_step}
    \x_{t_{i-1} \to t_i}  = \frac{\alpha_{t_i}}{\alpha_{t_{i-1}}} \tilde\x_{t_{i-1}} - \alpha_{t_i} \int_{\lambda_{t_{i-1}}}^{\lambda_{t_i}} e^{-\lambda} \hat\epsilonv_\theta(\hat\x_{\lambda}, \lambda)\dd\lambda.
\end{equation}
Therefore, to compute the value $\tilde\x_{t_i}$ for approximating $\x_{t_{i-1} \to t_i} $, we need to approximate the exponentially weighted integral of $\hat\epsilonv_\theta$ from $\lambda_{t_{i-1}}$ to $\lambda_{t_i}$.
~Denote $h_i\coloneqq \lambda_{t_i} - \lambda_{t_{i-1}}$, and $\hat\epsilonv_\theta^{(n)}(\hat\x_\lambda,\lambda)\coloneqq \frac{\dd^n \hat\epsilonv_\theta(\hat\x_\lambda,\lambda)}{\dd\lambda^n}$ as the $n$-th order total derivative of $\hat\epsilonv_\theta(\hat\x_\lambda, \lambda)$ w.r.t. $\lambda$. For $k\geq 1$, the $(k-1)$-th order Taylor expansion of $\hat\epsilonv_\theta(\hat\x_{\lambda},\lambda)$ w.r.t. $\lambda$ at $\lambda_{t_{i-1}}$ is
\begin{equation*}
    \hat\epsilonv_\theta(\hat\x_\lambda, \lambda) = \sum_{n=0}^{k-1} \frac{(\lambda - \lambda_{t_{i-1}})^n}{n!} \hat\epsilonv_\theta^{(n)}(\hat\x_{\lambda_{t_{i-1}}},\lambda_{t_{i-1}})+\Oc((\lambda - \lambda_{t_{i-1}})^{k}),
\end{equation*}
Substituting the above Taylor expansion into Eq.~\eqref{eqn:analytic_solution_each_step} yields
\begin{equation}
    \label{eqn:k-th-expansion}
    \x_{t_{i-1} \to t_i}\! =\! \frac{\alpha_{t_i}}{\alpha_{t_{i-1}}} \tilde\x_{t_{i-1}} - \alpha_{t_i} \sum_{n=0}^{k-1} \hat\epsilonv_\theta^{(n)}(\hat\x_{\lambda_{t_{i-1}}},\lambda_{t_{i-1}}) \!\int_{\lambda_{t_{i-1}}}^{\lambda_{t_i}}\!\! e^{-\lambda}  \frac{(\lambda - \lambda_{t_{i-1}})^n}{n!} \dd\lambda + \Oc(h_i^{k+1}),
\end{equation}
where the integral $\int e^{-\lambda}  \frac{(\lambda - \lambda_{t_{i-1}})^n}{n!} \dd\lambda$ can be \textbf{analytically} computed by repeatedly applying $n$ times of integration-by-parts (see Appendix~\ref{appendix:general_expansion}). Therefore, to approximate $\x_{t_{i-1} \to t_i}$, we only need to approximate the $n$-th order total derivatives $\hat\epsilonv_\theta^{(n)}(\hat\x_\lambda,\lambda)$ for $n\leq k-1$, which is a well-studied problem in the ODE literature~\citep{hochbruck2005explicit,luan2021efficient}.
By dropping the $\Oc(h_i^{k+1})$  error term and approximating the first $(k-1)$-th total derivatives with the ``stiff order conditions''~\citep{hochbruck2005explicit,luan2021efficient},
we can derive $k$-th-order ODE solvers for diffusion ODEs. We name such solvers as \textit{DPM-Solver} overall, and \textit{DPM-Solver-$k$} for a specific order $k$. %
Here we take $k=1$ for demonstration. In this case, Eq.~\eqref{eqn:k-th-expansion} becomes
\begin{align}
    \x_{t_{i-1} \to t_i} &= \frac{\alpha_{t_i}}{\alpha_{t_{i-1}}} \tilde\x_{t_{i-1}} - \alpha_{t_i} \epsilonv_\theta(\tilde\x_{t_{i-1}},t_{i-1}) \int_{\lambda_{t_{i-1}}}^{\lambda_{t_i}} e^{-\lambda}\dd\lambda + \Oc(h_i^2) \nonumber \\
    &= \frac{\alpha_{t_i}}{\alpha_{t_{i-1}}} \tilde\x_{t_{i-1}} - \sigma_{t_i} (e^{h_i} - 1)\epsilonv_\theta(\tilde\x_{t_{i-1}},t_{i-1}) + \Oc(h_i^2).   \nonumber
\end{align}
By dropping the high-order error term $\Oc(h_i^2)$, we can obtain an approximation for $\x_{t_{i-1} \to t_i} $. As $k=1$ here, we call this solver \textit{DPM-Solver-1}, and the detailed algorithm is as following.

\textbf{DPM-Solver-1.}\quad Given an initial value $\x_T$ and $M+1$ time steps $\{t_i\}_{i=0}^M$ decreasing from $t_0=T$ to $t_M=0$. Starting with $\tilde\x_{t_0}=\x_T$, the sequence $\{\tilde\x_{t_i}\}_{i=1}^M$ is computed iteratively as follows:
\begin{equation}
\label{eqn:1st}
    \tilde\x_{t_i} = \frac{\alpha_{t_i}}{\alpha_{t_{i-1}}} \tilde\x_{t_{i-1}} - \sigma_{t_i} (e^{h_i} - 1)\epsilonv_\theta(\tilde\x_{t_{i-1}},t_{i-1}),~~~~\mbox{where }h_i=\lambda_{t_i} - \lambda_{t_{i-1}}.
\end{equation}

For $k\geq 2$, approximating the first $k$ terms of the Taylor expansion needs additional intermediate points between $t$ and $s$~\citep{hochbruck2005explicit}.
The derivation is more technical so we defer it to Appendix~\ref{appendix:proof}.
Below we propose algorithms for $k=2,3$ and name them as \textit{DPM-Solver-2} and \textit{DPM-Solver-3}, respectively.










\begin{algorithm}[H]
    \centering
    \caption{DPM-Solver-2.}\label{alg:dpm-solver-2}
    \begin{algorithmic}[1]
    \Require initial value $\x_T$, time steps $\{t_i\}_{i=0}^M$, model $\epsilonv_\theta$
        \State $\tilde\x_{t_0}\gets\x_T$
        \For{$i\gets 1$ to $M$}
        \State $s_{i} \gets t_{\lambda}\left(\frac{\lambda_{t_{i-1}} + \lambda_{t_i}}{2}\right)$
        \State 
    $\uv_{i} \gets \frac{\alpha_{s_i}}{\alpha_{t_{i-1}}} \tilde\x_{t_{i-1}} - \sigma_{s_i}\left(e^{\frac{h_i}{2}} - 1\right)\epsilonv_\theta(\tilde \x_{t_{i-1}},t_{i-1})$
    \State  $\tilde\x_{t_{i}} \gets \frac{\alpha_{t_{i}}}{\alpha_{t_{i-1}}} \tilde\x_{t_{i-1}} - \sigma_{t_{i}}\left(e^{h_i} - 1\right)\epsilonv_\theta(\uv_{i},s_i)$
        \EndFor
        \State \Return $\tilde\x_{t_M}$
    \end{algorithmic}
\end{algorithm}
\begin{algorithm}[H]
    \centering
    \caption{DPM-Solver-3.}\label{alg:dpm-solver-3}
    \begin{algorithmic}[1]
    \Require initial value $\x_T$, time steps $\{t_i\}_{i=0}^M$, model $\epsilonv_\theta$
    \State $\tilde\x_{t_0}\gets\x_T$, $r_1\gets\frac{1}{3}$, $r_2\gets\frac{2}{3}$
    \For{$i\gets 1$ to $M$}
\State 
$s_{2i-1} \gets t_\lambda\left(\lambda_{t_{i-1}} + r_1 h_i\right),\quad s_{2i} \gets t_\lambda\left(\lambda_{t_{i-1}} + r_2 h_i\right)$ 
\State 
    $\uv_{2i-1} \gets \frac{\alpha_{s_{2i-1}}}{\alpha_{t_{i-1}}} \tilde\x_{t_{i-1}} - \sigma_{s_{2i-1}}\left(e^{r_1 h_i} - 1\right)\epsilonv_\theta(\tilde \x_{t_{i-1}},t_{i-1})$
\State 
    $\Dv_{2i-1} \gets \epsilonv_\theta(\uv_{2i-1},s_{2i-1}) - \epsilonv_\theta(\tilde\x_{t_{i-1}},t_{i-1})$
\State
    $\uv_{2i} \gets \frac{\alpha_{s_{2i}}}{\alpha_{t_{i-1}}} \tilde\x_{t_{i-1}} 
    - \sigma_{s_{2i}}\left(e^{r_2 h_i} - 1\right)\epsilonv_\theta(\tilde\x_{t_{i-1}},t_{i-1}) - \frac{\sigma_{s_{2i}}r_2}{r_1}\left( \frac{e^{r_2 h_i} - 1}{r_2 h_i} - 1\right)\Dv_{2i-1}$
\State 
    $\Dv_{2i} \gets \epsilonv_\theta(\uv_{2i},s_{2i}) - \epsilonv_\theta(\tilde\x_{t_{i-1}},t_{i-1})$
\State $\tilde\x_{t_{i}} \gets \frac{\alpha_{t_{i}}}{\alpha_{t_{i-1}}} \tilde\x_{t_{i-1}}
        - \sigma_{t_{i}}\left(e^{h_i} - 1\right)\epsilonv_\theta(\tilde\x_{t_{i-1}},t_{i-1}) - \frac{\sigma_{t_{i}}}{r_2}\left( \frac{e^{h_i} - 1}{h} - 1\right)\Dv_{2i}$
        \EndFor
        \State  \Return $\tilde\x_{t_M}$
    \end{algorithmic}
\end{algorithm}



Here, $t_\lambda(\cdot)$ is the inverse function of $\lambda(t)$, which has an analytical formulation for the practical noise schedule used in~\citep{ho2020denoising,nichol2021improved}, as shown in Appendix~\ref{appendix:implementation}. The chosen intermediate points are ($s_i$, $\uv_{i}$) for DPM-Solver-2 and $(s_{2i-1}, \uv_{2i-1})$ and $(s_{2i}, \uv_{2i})$ for DPM-Solver-3. 
As shown in the algorithm, DPM-Solver-$k$ requires $k$ function evaluations per step for $k=1, 2, 3$. Despite the more expensive steps, higher-order solvers ($k=2, 3$) are usually more efficient since they require much fewer steps to converge, due to their higher convergence order. 
We show that DPM-Solver-$k$ is $k$-th-order solver, as stated in the following theorem. The proof is in Appendix~\ref{appendix:proof}.



\begin{theorem}[DPM-Solver-$k$ as a $k$-th-order solver]
\label{thrm:order}
Assume $\epsilonv_\theta(\x_t,t)$ follows the regularity conditions detailed in Appendix~\ref{appendix:assumptions}, then for $k=1,2,3$, DPM-Solver-$k$ is a $k$-th order solver for diffusion ODEs, i.e., for the sequence $\{\tilde\x_{t_i}\}_{i=1}^M$ computed by DPM-Solver-$k$, the approximation error at time $0$ satisfies $\tilde \x_{t_M} - \x_0 =\Oc(h_{\max}^k)$, where $h_{\text{max}}=\max_{1\leq i \leq M}(\lambda_{t_i}-\lambda_{t_{i-1}})$.
\end{theorem}

Finally, solvers with $k\geq 4$ need much more intermediate points as shown by previous work~\citep{hochbruck2005explicit,luan2021efficient} for  exponential integrators. Therefore, we only consider $k$ from $1$ to $3$ in this work, while leaving the solvers with higher $k$ for future study.






























\subsection{Step Size Schedule}
\label{sec:adaptive}
The proposed solvers in Sec.~\ref{sec:high-order} need to specify the time steps $\{t_i\}_{i=0}^M$ in advance. We propose two choices of the time step schedule. One choice is handcrafted, which is to uniformly split the interval $[\lambda_T$, $\lambda_0$], i.e.
$\lambda_{t_i} = \lambda_T + \frac{i}{M}(\lambda_0 - \lambda_T)$, $i=0,\dots,M$. 
Note that this is different from previous work~\citep{ho2020denoising,song2020score} which chooses uniform steps for $t_i$.
Empirically, DPM-Solver with uniform time steps $\lambda_{t_i}$ can already generate quite good samples in few steps, where results are listed in Appendix~\ref{appendix:experiment}.
As the other choice, we propose an adaptive step size algorithm, which dynamically adjusts the step size by combining different orders of DPM-Solver. The adaptive algorithm is inspired by~\citep{jolicoeur2021gotta} and we defer its implementation details to Appendix~\ref{appendix:algorithm}.

For few-step sampling, we need to use up all the number of function evaluations (NFE). When the NFE is not divisible by $3$, we firstly apply DPM-Solver-3 as much as possible, and then add a single step of DPM-Solver-1 or DPM-Solver-2 (dependent on the reminder of $K$ divided by $3$), as detailed in Appendix~\ref{appendix:implementation}. In the subsequent experiments, we use such combination of solvers with the uniform step size schedule for NFE $\leq 20$, and otherwise the adaptive step size schedule.









\subsection{Sampling from Discrete-Time DPMs}
\label{sec:discrete-time}
Discrete-time DPMs~\citep{ho2020denoising} train the noise prediction model at $N$ fixed time steps $\{t_n\}_{n=1}^N$, and the noise prediction model is parameterized by $\tilde\epsilonv_\theta(\x_n, n)$ for $n=0,\dots,N-1$, where each $\x_n$ is corresponding to the value at time $t_{n+1}$. We can transform the discrete-time noise prediction model to the continuous version by letting
$\epsilonv_\theta(\x, t)\coloneqq \tilde\epsilonv_\theta(\x, \frac{(N-1)t}{T})$, for all $\x\in\R^d, t\in[0,T]$.
Note that the input time of $\tilde\epsilonv_\theta$ may not be integers, but we find that the noise prediction model can still work well, and we hypothesize that it is because of the smooth time embeddings (e.g., position embeddings~\citep{ho2020denoising}). By such reparameterization, the noise prediction model can adopt the continuous-time steps as input, and thus we can also use DPM-Solver for fast sampling.